% sec-02: the consortium structure.  Figure 1 is a d2-type nested container
% diagram of who contributes what and who receives what.
\section{Who Contributes What}

\begin{appfloat}[!tb]
\begin{appfig}[d2-type / nested consortium containers]
% Three containers on a 4.55cm pitch, three tiles each, plus a shared output
% container below.  Container titles sit above their children, never inside.
\foreach \i/\x/\a/\b/\c in {1/0/{Published protocol and IND}/{Audited simulation package}/{Analysis plan and code},
                            2/4.6/{Patients and operative expertise}/{Pharmacy, IRB, and monitoring}/{Pathology and specimen handling},
                            3/9.2/{Investigational drug supply}/{Regulatory cross-reference}/{Assay and bioanalytical support}}{
  \node[d2soft,text width=36mm,minimum height=8mm,anchor=north west] (x\i a) at (\x,0) {\a};
  \node[d2soft,text width=36mm,minimum height=8mm,anchor=north west] (x\i b) at ([yshift=-2mm]x\i a.south west) {\b};
  \node[d2soft,text width=36mm,minimum height=8mm,anchor=north west] (x\i c) at ([yshift=-2mm]x\i b.south west) {\c};}
\begin{scope}[on background layer]
\node[d2cont,fit=(x1a)(x1c),inner sep=6pt] (c1) {};
\node[d2cont,fit=(x2a)(x2c),inner sep=6pt] (c2) {};
\node[d2cont2,fit=(x3a)(x3c),inner sep=6pt] (c3) {};
\end{scope}
\node[d2title,anchor=south west] at ([yshift=1.2mm]c1.north west) {Independent scientist};
\node[d2title,anchor=south west] at ([yshift=1.2mm]c2.north west) {Academic cancer center};
\node[d2title2,anchor=south west] at ([yshift=1.2mm]c3.north west) {Drug developer};
\node[d2key,text width=112mm,minimum height=9mm,anchor=north] (out) at ([yshift=-9mm]c2.south)
  {Shared output: a specimen-level perioperative pharmacodynamic dataset, released openly, owned by none of the three};
\draw[d2edgeb] (c1.south) -- (out.north west);
\draw[d2edgeb] (c2.south) -- (out.north);
\draw[d2edgeb] (c3.south) -- (out.north east);
\end{appfig}
\vspace{-0.7cm}
\figcaption{The three contributing parties and the single shared output. Each
container holds only what that party can supply, and the output belongs to none
of them, which is what makes the question pre-competitive.}
\end{appfloat}

The order in which the three parties commit matters more than the totals. A
drug cross-reference that arrives after the IRB submission forces a resubmission;
a site agreement that arrives after the drug supply wastes the supply window.

\begin{appfloat}[!tb]
\begin{appfig}[mermaid-type / sequence diagram]
% Four lifelines 3.05 apart; eight messages 0.72 apart, so no arrow label
% touches the message above or below it.
\foreach \i/\x/\lab in {1/0/{Independent scientist}, 2/3.1/{FNIH}, 3/6.2/{Drug developer}, 4/9.3/{Academic center}}{
  \node[mmactor] (a\i) at (\x,0) {\lab};
  \draw[mmlife] (\x,-0.42) -- (\x,-6.35);}
\foreach \y/\f/\t/\msg in {%
  -0.85/0/3.1/{1. Concept and published protocol},
  -1.57/3.1/6.2/{2. Pre-competitive question confirmed},
  -2.29/6.2/3.1/{3. Cross-reference letter issued},
  -3.01/3.1/9.3/{4. Site feasibility request},
  -3.73/9.3/3.1/{5. Investigator and accrual estimate},
  -4.45/3.1/0/{6. Consortium agreement executed},
  -5.17/0/9.3/{7. IRB submission with fixed analysis plan},
  -5.89/9.3/0/{8. First participant screened}}{
  \draw[mmmsg] (\f,\y) -- (\t,\y) node[midway,above,font=\tiny\sffamily,fill=protowhite,inner sep=1.2pt] {\msg};}
\node[font=\tiny\itshape,text=protogray,anchor=north west,text width=100mm]
  at (-0.6,-6.55) {Message 3 must precede message 7. A cross-reference that
  arrives after the IRB submission forces a resubmission and costs the accrual
  window.};
\end{appfig}
\vspace{-0.7cm}
\figcaption{The order in which the four parties must act. The dependency that
governs the schedule is message 3 preceding message 7, not the size of any one
contribution.}
\end{appfloat}
